% from main.tex
%%%%%%%%%%  scan idx 524  |  folio 511  %%%%%%%%%%
homomorphisme se factorise en :
\[
K_\bullet(S)\otimes_{K^\circ(S)}K^\circ(X)\otimes_{K^\circ(X)}K^\circ(D(E))
\longrightarrow K_\bullet(X)\otimes_{K^\circ(X)}K^\circ(D(E))
\longrightarrow K_\bullet(D(E))
\]
et le premier homomorphisme est injectif. Comme, d'après VI 4.3, 4.6 et 4.7, $K^\circ(D(E))$ est libre sur $K^\circ(X)$, cela entraîne l'injectivité de :
\[
K_\bullet(S)\otimes_{K^\circ(S)}K^\circ(X)\to K_\bullet(X).
\]
Soit donc $X=D(E)$. On procède par réccurrence sur le rang de $E$. Le cas $n=1$ étant trivial, supposons la propriété prouvée pour $n-1$. On introduit alors $P=\mathbb P(E)$ et le noyau $F$ de l'homomorphisme canonique $p^*(E)\to \mathcal O_P(1)$, où $p$ est le morphisme structural de $P$. Alors $X$ s'identifie au schéma des drapeaux de $F$ sur $P$, et on voit d'après l'hypothèse de récurrence qu'il suffit de prouver 3.1 pour $P$. Compte tenu de VI 1.1, le résultat à démontrer peut donc s'énoncer :

\textbf{Corollaire 3.2.} Soient $S$ un schéma noethérien, $E$ un $\mathcal O_S$-Module localement libre de rang $n$, $X=\mathbb P(E)$ le fibré projectif associé, $f$ le morphisme structural de $X$, et $\xi$ la classe dans $K^\circ(X)$ du faisceau $\mathcal O_X(1)$. Alors l'homomorphisme :
\[
\varphi:(K_\bullet(S))^n\longrightarrow K_\bullet(X)
\]
défini par :
\[
\varphi((x_i))=\sum_{i=0}^{n-1}\xi^{-i}f^*(x_i)
\]
est un isomorphisme.

On a déjà vu que cet homomorphisme est surjectif, puisque les $\xi^i$, $i=0,\ldots,-n+1$, forment une base de $K^\circ(X)$ sur $K^\circ(S)$; pour voir qu'il est injectif, on utilise l'homomorphisme $f_*$, défini car $f$ est projectif. On a
